\documentclass[10pt]{article}

\usepackage[utf8]{inputenc}

\usepackage[T1]{fontenc}

\usepackage{amsmath}

\usepackage{amsfonts}

\usepackage{amssymb}

\usepackage[version=4]{mhchem}

\usepackage{stmaryrd}



\begin{document}

сов вычетов по модулю $m$ имеет наиболее простую структуру, являясь циклическоӥ группой порядка $\varphi(m)$. В остальных случаях групшы приведенных классов вычетов устроены значительно сложнее.



Многие задачи теории чисел сводятся к вопросу о разрешимости или неразрешимости того или иного типа С. Ввиду этого теория С., впервые систематически изложенная К. Гауссом (см. [5]) и поставленная им в основу классической теории чисел, до сих пор является одним из основных средств при решении теоретико-числовых задач. В этой связи фундаментальное значение для теории чисел имеет исследование вопроса о числе решений алгебраич. С. Простейшим типом алгебраич. С. являются С. 1-й стешени с одним неизвестным $a x \equiv b(\bmod m)$. Полный ответ на вопрос о числе решений С. 1-й степени дает следующая теорема. Пусть $(a, m)=d$. Сравнение $a x \equiv b(\bmod m)$ неразрешимо, если $b$ не делится на $d$. При $b$, кратном $d$, С. имеет ровно $d$ решений.



Вопрос о разрешимости системы линейных С.



\[

\sum_{j=1}^{S} a_{i j} x_{j} \equiv b_{i}(\bmod p), i=1,2, \ldots, t,

\]



по простому модулю $p>2$ полностью решается в общей теории линейных уравнений над произвольными полями. Случай составного модуля можно свести к случаю простого модуля.



Следующим по сложности изучения видом алгебраич. С. с одним неизвестным является двучленное сравнение



\[

x^{n} \equiv a(\bmod m) \text { при }(a, m)=1 .

\]



Если С. $x^{n} \equiv a(\bmod m)$ имеет решения, то $a$ наз. в ыч е то м с т е п е н и $n$ по мод улю тивном случае $a$ наз. не в ы ч е то м с т е п е ни $n$ п о м о д у л ю $m$. В частности, при $n=2$ вычеты или невычеты называются к в а д р а т и ч ны м и, при $n=3-$ к у б и ч е с кими, а при $n=4-6$ и к в а др а т ичны ми.



Вотрос о числе решений $\mathrm{C}$.



\[

f(x) \equiv 0(\bmod m)

\]



по составному модулю $m=p_{1}^{\alpha_{1}} \ldots p_{s}^{\alpha_{s}}$ сводится к вопросу о числе решений $\mathrm{C}$.



\[

f(x) \equiv 0\left(\bmod p_{i}^{\alpha_{i}}\right)

\]



для каждого $i=1,2, \ldots, s$. Имеет место следующая теорема: если $m_{1}, \ldots, m_{r}$ - попарно взаимно простые целые положительные числа, то число $N$ решениі̆ $\mathrm{C}$. $f(x) \equiv 0\left(\bmod m_{1} \ldots m_{r}\right)$ выражается через величины $N_{i}$, $i=1,2, \ldots, r$, равные числу решений соответств ующих С. $f(x) \equiv 0\left(\bmod m_{i}\right)$ по формуле $N=\prod_{i=1}^{r} N_{i}$. В свою очередь, вопрос о числе решений $\mathrm{C}$.



\[

f(x) \equiv 0\left(\bmod p^{\alpha}\right), \alpha>1,

\]



в основном сводится к вопросу о числе решений $\mathrm{C}$. $f(x) \equiv 0(\bmod p)$ по простому модулю $p$. Основой для такого сведения служит следующая теорема: всякое решение $x \equiv x_{1}(\bmod p)$ сравнения $f(x) \equiv 0(\bmod p)$ единственным образом продолжается до решения $x \equiv x_{\alpha}\left(\bmod p^{\alpha}\right)$ сравнения $f(x) \equiv 0\left(\bmod p^{\alpha}\right)$, если только формальная производная $f^{\prime}\left(x_{1}\right) \not 0(\bmod p)$, т. е. найдется единственное реншение $x_{\alpha}$ С. $f(x) \equiv 0\left(\bmod p^{\alpha}\right)$ такое, что $x_{\alpha}=x_{1}(\bmod p)$.



Аналогичная ситуация имеет место и в случае алгебраич. С. от нескольких переменных, т. е. в невырожденных случаях вопрос о числе решений С. $F\left(x_{1}\right.$, $\left.\ldots, x_{n}\right)=0(\bmod m)$ по составному модулю $m$ редуцируется к вопросу о числе рептений такого же С. по простому модулю. Верхнюю границу для числа решений С. $f(x) \equiv$ $\equiv 0(\bmod p), \quad$ где $f(x)=a_{0} x^{n}+a_{1} x^{n-1}+\ldots+a_{n}, \quad a_{0} \neq$ $\not 0(\bmod p)$, дает т е о р м а ЈI а г р а н ж а: число решений $\mathrm{C}$. $f(x) \equiv 0(\bmod p)$ не превосходит степени многочлена $f(x)$. Изучение вопроса о точном числе решениї С. $f(x) \equiv 0(\bmod p)$ в общем случае встречает значительные трудности, и в атом направлении к настоящему времени (1984) получены липь немногие результаты.



Утверждение теоремы Лагранжа перестает быть верным для составного модуля. Указанное специфическое свойство простых чисел объясняется тем, что классы вычетов по простому модулю $p$ образуют поле (см. Сравнение по простолу модулю). Другое специфическое свойство простых чисел, объясняемое тем же фактом, выражает Вильсона теорема.



При изучении квадратичных $\mathrm{C}$.



\[

x^{2} \equiv a(\bmod p)

\]



по простому модулю $p$ и для их приложений вводятся Лежандра символ и Якоби символ.



Алгебраич. С. по простому модулю с двумя неизвестными (и вообще с любым числом неизвестных)



\[

F(x, y) \equiv 0(\bmod p)

\]



можно трактовать как уравнения над простым конечным полем из $p$ элементов. Поэтому для их изучения наряду с методами теории чисел применяются методы теории алгебраич. функций и алгебраич. геометрии. Именно этими методами была впервые получена асимптотика вида



\[

N_{p}=p+O(\sqrt{p})

\]



для числа $N_{p}$ решений широкого класса С. $F(x, y) \equiv$ $\equiv 0(\bmod p)$. Такая же асимптотика была впоследствии выведена методами теории чисел без выхода за рамки теории С. (см. [6]).



Об алгебраич. С. от многих неизвестных известно сравнительно мало. В качестве общего результата можно привести следующее (т е о р е м а III е в а лл е). Пусть $F\left(x_{1}, \ldots, x_{n}\right)$ - многочлен с целыми рациональными коэффициентами, степень к-рого меньше n. Тогда число решений С.



\[

F\left(x_{1}, \ldots, x_{n}\right) \equiv 0(\bmod p)

\]



делится на простое тисло $p$. Сильныӥ результат в этом круге вопросов получен ПІ. Делинем [9] (см. также $[10])$.



Большое значение имеет также исследование вопроса о числе решений С. на неполной системе вычетов. Систематич. решение этого вопроса было начато работами И. М. Виноградова (см. [4]). В качестве примера задач такого типа может служить задача о распределении квадратичных вычетов и невычетов во множестве $1,2, \ldots, p-1$, связанная с изучением решений $\mathrm{C}$. $y^{2} \equiv x(\bmod p)$ (см. Виноградова гипотезь, Распределение степенных вычетов и невычетов).



Лum.: [1] Б о p е в ич Э. И., III а ф а р е в и ч И. Р., Теория чисел, 2 изд., М., 1972, [2] в е н к о в Б. А., Элементарноя теория чисел, М. 8 изд., М., 1972; [4] В и н о г р ад о в И. М., Избр. труды, М., 1952; [5] Г а у с с ІК. Ф., Труды по теории чисел, пер. с нем., М., 1959; [6] С т е п а н о в С. А.,

"Тр. Матем. ин-та АН СССР", 1973, т. 132, с. $237-46 ;$; [7] Х а с"Тр. Матем. ин-та АН СССР", 1973, т. 132, с. 237-46; [7] Х а сд р а с е к х а р а н К., Введение в аналитическую теорию чиIHES", 1973, v. 43, p. 273-307; [10] K a t z N. M., "Proc. Simp. Pure Math.", 1976, v. 28, p. 275-305. "C. A. Степанов.



СРАВНЕНИЕ ОТ НЕСКОंЈЫКИХ ПЕРЕМЕННЫХ сравнение вида



\[

f\left(x_{1}, \ldots, x_{n}\right) \equiv 0(\bmod m),

\]



где $f\left(x_{1}, \ldots, x_{n}\right)$ - многочлен от $n \geqslant 2$ переменных с целыми рациональными коәффициентами, не все из





\end{document}